\documentclass[book.tex]{subfiles}

\begin{document}

\chapter{Energy-based-models}
\label{chap:energy_models}
\noindent\rule{11cm}{0.4pt} \\
\textbf{Prerequisites:}  \autoref{chap:ml_basics}, \autoref{chap:grad_descent} \\
\textbf{Difficulty Level:} ** \\
\noindent\rule{11cm}{0.4pt}
\vspace{5mm}

Energy based models adapt ideas from statistical mechanics to create a tool for modeling probability distributions. We will return to study the relevant statistical mechanics in more depth later in the book, but in this chapter will focus on their capacity as machine learning models. 

\section{Boltzmann Machine}

The Boltzmann Machine is a form of an energy based model. An example of this is given in Figure \ref{fig:boltzmann_machine}. The energy function for a BM is given by
\begin{equation}
E(v,h)= -b\sum_{i=1}^{n_v} v_i -c\sum_{i=1}^{n_h} h_i 
- \sum_{i=1}^{n_v}\sum_{j=1}^{n_v} v_i v_j u_{i,j}
- \sum_{i=1}^{n_h}\sum_{j=1}^{n_h} h_i h_j v_{i,j}
- \sum_{i=1}^{n_v}\sum_{j=1}^{n_h} v_i h_j w_{i,j}
\end{equation}

Here $v$ is the visible layer, and $h$ is the hidden layer. The statistical mechanics experts reading along may note the similarity to the Hamiltonian for the Ising model.

These large networks often take a long time to reach their equilibrium distributions, when distributions are complex multi-modal distribution. These difficulties can be overcome by simplifying the form of the energy function as we can see in the next section.  

\begin{marginfigure}
% \includegraphics[width=\textwidth]{figures/Differentiable Models/autoencoders/Boltzmannexamplev1.png}
  \includegraphics[width=\textwidth]{figures/Differentiable Models/autoencoders/Boltzmann machine (2).png}
\caption{Example of a Boltzmann Machine}
\label{fig:boltzmann_machine}
\end{marginfigure}

\begin{lstlisting}[language=python]
class BoltzmannMachine(b: $\mathbb{R}$, c: $\mathbb{R}$, U: $\mathbb{R}^{n_v \times n_v}$, V: $\mathbb{R}^{n_h \times n_h}$, W: $\mathbb{R}^{n_h \times n_v}$):

  def $\lambda$(v: $\mathbb{R}^{n_v}$, h: $\mathbb{R}^{n_h}$) $\to \mathbb{R}$ :
    return (-b * sum(v) - c * sum(h) - sum(U * (v $\otimes$ v)) 
            - sum(V * (h $\otimes$ h)) - sum(W * (v $\otimes$ h)))

\end{lstlisting}

\subsection{Restricted Boltzmann Machines}
Restricted Boltzmann machines are some of the most common building blocks of deep probabilistic models. RBMs are undirected probabilistic graphical models containing a layer of observable variables and a single layer of latent variables. RBMs may be stacked (one on top of the other) to form deeper models. It is a bipartite graph, with no connections permitted between any variables in the observed layer or between any units in the latent layer (see Fig \ref{fig:rbm}).

\begin{marginfigure}
% \includegraphics[width=\textwidth]{figures/Differentiable Models/autoencoders/restricted_boltzmann_machine.png}
\includegraphics[width=\textwidth]{figures/Differentiable Models/autoencoders/Restricted Boltzmann machine1.png}

\caption{Example of a Restricted Boltzmann Machine}
\label{fig:rbm}
\end{marginfigure}

Let the observed layer consist of a set of $n_v$ binary random variables which will be referred to collectively with the vector $v$. The latent or hidden layer of the network will consist of a set of $n_h$ binary random variables and will be referred to as $h$. The restricted Boltzmann machine is also an energy based model with the joint probability distribution specified by its energy function:
\begin{align}
P(v, h)=\frac{1}{Z}\exp(-E(v,h))
\end{align}
The energy function for an RBM (in close analogy to the mean field approximation of the Ising Model) is given by
\begin{align}
E(v,h)= -b\sum_{i=1}^{n_v} v_i -c\sum_{i=1}^{n_h} h_i - \sum_{i=1}^{n_v}\sum_{j=1}^{n_h} v_i h_j w_{i,j}
\end{align}
and the partition function, $Z$, is given by:
\begin{align}
Z=\sum_v \sum_h \exp\left(-E(v,h)\right)
\end{align}
It is apparent from the definition of the partition function $Z$ that the naive method of computing $Z$ (exhaustively summing over all states) could be computationally intractable, unless a cleverly designed algorithm could exploit regularities in the probability distribution to compute $Z$ faster. In the case of restricted Boltzmann machines, the partition function $Z$ is intractable. The intractable partition function $Z$ implies that the normalized joint probability distribution $P(v)$ is also intractable to evaluate.

Though $P$ is intractable, the bipartite graph structure of the RBM has the very special property that its conditional distributions $P(h|v)$ and $P(v|h)$ are factorial and relatively simple to compute and to sample from. Deriving the conditional distributions from the joint distribution is straightforward:
\begin{align}
P(h|v) &= \frac{P(h,v)}{P(v)} = \frac{1}{Z'}\prod^{n_h}_{j=1}\exp\Big(c_jh_j+v^TW_{:,j}h_j\Big)
\end{align}
Since the equations are conditioning on the visible units v, they can be treated as constant with respect to the distribution $P(h|v)$. The factorial nature of the conditional $P(h|v)$ follows immediately from our ability to write the joint probability over the vector h as the product of (unnormalized) distributions over the individual elements, $h_j$. It is now a simple matter of normalizing the distributions over the individual binary $h_j$.
\begin{align}
P(h_j=1|v)=\sigma(c_j+v^TW_{:,j})
\end{align}
%The full conditional over the hidden layer as a factorial distribution can now be expressed as:
%\begin{align}
%P(h|v) &= \prod^{n_h}_{j=1}\sigma\Big((2h-1)(c+W^Tv)\Big)_j
%\end{align}
A similar derivation will show the following expression for the other condition of interest, $P(v|h)$:
\begin{align}
P(v_j = 1|h)=\sigma\Big(b_i+W_{i,:}h\Big)
    \label{visible}
\end{align}

\begin{lstlisting}[language=python]
class RestrictedBoltzmannMachine(b: $\mathbb{R}$, c: $\mathbb{R}$, W: $\mathbb{R}^{n_h \times n_v}$):

  def $\lambda$(v: $\mathbb{R}^{n_v}$, h: $\mathbb{R}^{n_h}$) $\to \mathbb{R}$:
    return (-b * sum(v) - c * sum(h) - sum(W * (v $\otimes$ h)))
\end{lstlisting}

\subsection{Training Restricted Boltzmann Machines}
Because the RBM admits efficient evaluation and differentiation of $\Tilde{P}(v)$ and efficient MCMC sampling in the form of block Gibbs sampling, it can readily be trained with techniques such as Contrastive Divergence and Stochastic Maximum Likelihood. Compared to other undirected models used in deep learning, the RBM is relatively straightforward to train because $P(h|v)$ can be computed exactly in closed form. The update equations for this model are:
\begin{align}
    -\frac{\partial \log p(v)}{\partial W_{ij}} &= -v_j^{(i)}\sigma(W_iv^{(i)}+c_i)+E[p(h_i|v)v_j]
    \label{weights1} \\
    -\frac{\partial \log p(v)}{\partial c_i} &= -\sigma(W_iv^{(i)}+c_i)+E[p(h_i|v)]
    \label{weights2} \\
    -\frac{\partial \log p(v)}{\partial b_j} &= -v_j^{(i)}+E[p(v_j|h)]
    \label{weights3} \\
\end{align}

\subsection{Contrastive Divergence}
The gradient equation that needs to be solved with the partition function is:
\begin{align}
    \nabla_\theta \log Z &= E_{x\sim p(x)} \left [ \nabla_\theta \Tilde{p}(x) \right ]
\end{align}
The naive way of implementing this equation is to compute it by burning in a set of Markov chains from a random initialization every time the gradient is needed. When learning is performed using stochastic gradient descent, this means the chains must be burned in once per gradient step. The high cost of burning in the Markov chains in the inner loop makes this procedure computationally infeasible, but this procedure is the starting point that other more practical algorithms aim to approximate. 

This approach can viewed as trying to achieve balance between two forces, one pushing up on the model distribution where the data occurs, and another pushing down on the model distribution where the model samples occur. The two forces correspond to maximizing $\log \Tilde{p}$ and minimizing $\log Z$. Several approximations to the negative phase are possible. Each of these approximations can be understood as making the negative phase computationally cheaper but also making it push down in the wrong locations.
\begin{figure}
% \includegraphics[width=\textwidth]{figures/Differentiable Models/autoencoders/contrastive_divergence.png}
\label{cd}
\includegraphics[width=\textwidth]{figures/Differentiable Models/autoencoders/contrastive divergence1.png}
\label{cd}
\caption{Illustration of contrastive divergence.} %\textcolor{red}{TODO: Replace}}
\end{figure}

The contrastive divergence (CD, or CD-$k$ to indicate CD with $k$ Gibbs steps) algorithm initializes the Markov chain at each step with samples from the data distribution. Obtaining samples from the data distribution is free, because they are already available in the data set. Initially, the data distribution is not close to the model distribution, so the negative phase is not very accurate. Fortunately, the positive phase can still accurately increase the model\textquotesingle s probability of the data. After the phase has had some time to act, the model distribution is closer to the data distribution, and the negative phase starts to become accurate. 

The pseudo-code of the contrastive divergence algorithm is as follows:
\begin{enumerate}
    \item Set the visible units to a training vector
    \item While not converged:
    \begin{enumerate}
        \item Update all hidden units in parallel using Eq. \ref{hidden}
        \item Update all visible units in parallel to get ``reconstructions'' using Eq. \ref{visible}
        \item Update all hidden units in parallel again using Eq. \ref{hidden}
        \item Update the coefficients of the energy model using Eqs. \ref{weights1}-\ref{weights3}
        \item Select another training vector
    \end{enumerate}
\end{enumerate}

\noindent We can convert our description into Physika code

 \begin{lstlisting}[language=python]
# Set $\epsilon$, the step size, to a small positive number
# Set $k$, the number of Gibbs steps high enough to allow a 
# Markov chain sampling from $p(x;\theta)$ to mix when 
# initialized from $p_{data}$      
class BoltzmannMachine:
  def contrastive_divergence(X: $\mathbb{R}^{N\times d}$,  $\epsilon: \mathbb{R}$, $k: \mathbb{N}$):
    while True:      # While not converged:
      # Sample a minibatch of m examples 
      # $\{x^{(1)},..,x^{(m)}\}$ from the training set
      $\{x^{(1)},..,x^{(m)}\}$ = Sample(X)
      $g = \frac{1}{m}\sum_{i=1}^m \nabla_\theta \log\ \Tilde{p}(x^{(i)};\theta)$
      for i = 1 to m:
        $\Tilde{x}^{(i)}=x^{(i)}$
      for i = 1 to k 
        $\Tilde{h}^{(i)}=\textrm{gibbs_update}(\Tilde{h}^{(i)}, \Tilde{x})$
      for j = 1 to m 
        $\Tilde{x}^{(j)}=\textrm{gibbs_update}(\Tilde{x}^{(j)}, \Tilde{h})$
      for i = 1 to k 
        $\Tilde{h}^{(i)}=\textrm{gibbs_update}(\Tilde{h}^{(i)}, \Tilde{x})$
      $g = g - \frac{1}{m}\sum_{i=1}^m \nabla_\theta \log\ \Tilde{p}(\Tilde{x}^{(i)};\theta)$
      $\theta = \theta + \epsilon g$
 \end{lstlisting}
 
Of course, CD is still an approximation to the correct negative phase. The main way that CD qualitatively fails to implement the correct negative phase is that it fails to suppress regions of high probability that are far from actual training examples. These regions that have high probability under the model but low probability under the data generating distribution are called spurious modes. Essentially, it is because modes in the model distribution that are far from the data distribution will not be visited by Markov chains initialized at training points, unless k is very large. The CD estimator is biased for RBMs and fully visible Boltzmann machines, in that it converges to different points than the maximum likelihood estimator. Since the bias is small, CD could be used as an inexpensive way to initialize a model that could later be fine-tuned via more expensive MCMC methods.

\section{Performing Classification}

To use a Restricted Boltzmann machine to solve classification tasks, the visible vector is set to a concatenation of the feature vector and the label 
\begin{align}
v &= x_i | y_i
\end{align}
For test data, the visible vector is partially initialized with just the feature vector and the label is sampled from the distribution.

\end{document}